\begin{figure*}[t]
\centering
\begin{subfigure}{0.48\textwidth}
\includegraphics[width=0.95\linewidth]{figures/predictability_profile.pdf}
\caption{Predictability across depth}
\label{fig:momentum-profile}
\end{subfigure}
\hfill
\begin{subfigure}{0.48\textwidth}
\includegraphics[width=0.95\linewidth]{figures/recovery_vs_p.pdf}
\caption{Recovery vs.\ predictability}
\label{fig:momentum-recovery}
\end{subfigure}
\caption{\textbf{Transformer depth has predictable local momentum.} \emph{Left:}
held-out explained update energy $P_\ell$ (\cref{eq:pscore}) against normalized
layer depth, for each model; markers indicate the layers selected for skipping.
Predictability is concentrated in the middle of the network, and the profile has
the same qualitative shape at every scale. \emph{Right:} for each layer, the
fraction of single-layer plain-skip NLL damage that \mname recovers, plotted
against that layer's $P_\ell$. Layers whose updates are predictable are the layers
where update prediction repairs the most damage, so $P_\ell$ --- computable from
the dense model on unlabeled text, before any skipping decision --- forecasts where
skipping is safe. Measurements for 0.5B and 1.5B are made on an NVIDIA TITAN X
(Pascal) GPU and for Qwen2.5-\topscale on a single NVIDIA RTX 3090 GPU.}
\label{fig:momentum}
\end{figure*}
